% Profile: runware-creative -- Creative Engineer / AI Creative Technologist,
% Runware (France, remote; generative image and video inference platform;
% posting read 2026-10-07).
%
% What the posting asks for, and what this profile puts in front of it:
%
%   "Experiment hands-on with new generative image and video models"
%        -> the coding-agent bullet, first: GLM-5.3 with vision capabilities,
%           self-hosted and served with vLLM (from Hugging Face,
%           \cvshowstackdetailtrue). A vision model, not an image or video
%           generator -- the summary says so rather than blur it
%   "Develop repeatable and automated generation workflows ... at scale"
%        -> the VLM labeling bullet: automated vision-language pipelines that
%           cut hand-labeling by 75%
%   "Coding and API experience, including building custom integrations,
%    scripts, tools"
%        -> the summary's "automated pipelines and tools". The AI-gateway
%           bullet (custom Go plugins in front of the model APIs) is the
%           stronger evidence but did not fit the page
%   "A strong creative eye" / "technical depth"
%        -> the Projects section's Vulkan/OpenGL renderers (since 2020) and
%           the summary's B.Sc. thesis on blue-noise error distribution for
%           real-time image quality. Image quality is the owner's own
%           research topic, and the renderers are public
%   "Strong technical communication skills" / "turn complex experiments
%    into clear, engaging content"
%        -> the summary: four years of teaching, and the public
%           documentation pipeline (books, slides, this CV)
%   "practical experience across multiple image and/or video models"
%        -> SemProbe, the ECML PKDD 2026 demo (Publications), which uses
%           diffusion-based inpainting to probe object detectors. One
%           generative image technique used in research, named in the
%           summary; no model family is claimed beyond it
%   "Experience with model adaptation or fine-tuning, such as LoRAs"
%        -> the LoRA tag, which the experience section keeps because the
%           fine-tuning is real (an LLM, not a diffusion model)
%   "Contributions to open-source generative AI projects or communities"
%        -> not claimed as such; the summary names the owner's public
%           repositories and upstreamed fixes, which is open source but not
%           generative media
%   "Work independently and move quickly"
%        -> the coding agent, started on the owner's own initiative
%   France, remote -> no relocation line: the job is remote. French B1 in
%   section_private.
%
% Not claimed: ComfyUI or comparable node-based workflow tooling, custom
% ComfyUI nodes, any image or video generation model (FLUX, Stable Diffusion,
% Wan, LTX, Hunyuan, Kling), a generative-AI portfolio, social-media or
% launch content, LoRA training for diffusion models, work at a generative-AI
% or inference platform company, public talks or livestreams (the YouTube
% handle stays commented out in cv.tex: no videos published), and a
% three-years-in-generative-media figure. These are the posting's core
% requirements; the cover letter, if one is written, is where to say how the
% owner would close them, not this page.
%
% Nothing is deleted for this profile: every bullet it does not show is only
% switched off here and stays in section_experience.tex.

\newcommand{\cvTagline}{\IfLanguageName{english}{Research Associate @ Fraunhofer IPA | Vision models, GPU pipelines \& real-time rendering}{Wissenschaftlicher Mitarbeiter @ Fraunhofer IPA | Vision-Modelle, GPU-Pipelines \& Echtzeit-Rendering}}

\cvshowschoolfalse
\cvshowthesislinesfalse
% Off to fit one page: the gateway plugins (the posting's "custom
% integrations"), the self-labeling and MLflow/ClearML clauses, and the CAS
% Software bullet. The page ran about twelve lines over with them on.
\cvshowgatewayfalse
\cvshowmldetailfalse
\cvshowcasdetailfalse
% Hugging Face as the model source: the ecosystem this posting lives in.
\cvshowstackdetailtrue
% Not rendered here, kept in the source: the industrial computer-vision,
% build-system and HPC depth this reader needs least, traded for the
% Projects section.
\cvshowremanetfalse
\cvshowberghoffalse
\cvshowsafetyfalse
\cvshowkiprotfalse
\cvshowquantumfalse
\cvshowcudakernelsfalse
\cvshowbuildsfalse
\cvshowdocsfalse
% SemProbe's detail line stays off: the summary already says what the demo
% does, in fewer lines than the Publications column would spend on it.
\cvshowsemprobefalse
% Only the renderer stays. The blue-noise thesis is in the summary and in the
% B.Sc. tags (Path Tracing, Blue Noise), so its entry went to fit the page.
\cvshowprojbluenoisefalse
\cvshowprojmlfalse
\cvshowprojmrfalse
\cvshowprojcollisionfalse
\setlength{\cvtagsep}{3pt}
% The posting's ecosystem first (Python, PyTorch, Hugging Face, LoRA, vLLM),
% then the GPU and graphics side.
\renewcommand{\cvFhgTags}{Python, PyTorch, Hugging Face, LoRA, vLLM, CUDA, C++, Docker, uv, GitLab CI/CD}

% The vision model first, then the automated pipelines, then the GPU
% platform under them. Every bullet is listed, hidden ones too.
\renewcommand{\cvFraunhoferOrder}{%
	\cvFhgAgent\cvFhgVlm\cvFhgGateway\cvFhgCicd\cvFhgRemanet\cvFhgBerghof%
	\cvFhgSafety\cvFhgKiprot\cvFhgQuantum\cvFhgKernels\cvFhgBuilds\cvFhgDocs}

\newcommand{\cvBody}{%
    % Summary for the runware-creative profile: Creative Engineer / AI Creative
% Technologist, Runware (France, remote).
%
% Carries what the bullets cannot: that the owner's background is computer
% graphics and image quality (the Projects section is the evidence), that the
% model work below is vision, the one generative technique on record is
% SemProbe's diffusion-based inpainting, and the technical-communication evidence:
% four years of teaching and the public documentation pipeline. See the
% not-claimed list in profiles/runware-creative.tex before adding ComfyUI or
% any generation model here.
\IfLanguageName{english}{
    \par{I come from computer graphics: my B.Sc.\ thesis was on image quality in real-time rendering, and I have maintained my own Vulkan and OpenGL renderers since 2020. At Fraunhofer IPA I self-host and tune vision-language models on our GPUs and turn them into automated pipelines and tools that researchers use every day, and I co-authored an ECML PKDD demo that edits images with diffusion-based inpainting. I taught computer organization to undergraduates for four years, and I publish the documentation pipeline that typeset this CV, the books and the slides I write.}
}{
    \par{Ich komme aus der Computergrafik: Meine Bachelorarbeit behandelte Bildqualität im Echtzeit-Rendering, und seit 2020 pflege ich eigene Vulkan- und OpenGL-Renderer. Am Fraunhofer IPA betreibe und optimiere ich Vision-Language-Modelle auf unseren GPUs und mache daraus automatisierte Pipelines und Werkzeuge, die Forschende täglich nutzen; außerdem bin ich Co-Autor einer ECML-PKDD-Demo, die Bilder per Diffusions-Inpainting verändert. Ich habe vier Jahre Rechnerorganisation gelehrt und veröffentliche die Dokumentationspipeline, die diesen Lebenslauf, meine Bücher und Folien setzt.}
}
%
    \vspace{0.25cm}%
    % YAAC Another Awesome CV LaTeX Template
%
% This template has been downloaded from:
% https://github.com/darwiin/yaac-another-awesome-cv
%
% Author:
% Christophe Roger
%
% Template license:
% CC BY-SA 4.0 (https://creativecommons.org/licenses/by-sa/4.0/)

%Section: Scholarships and additional info
\IfLanguageName{english}{
    
    \sectionTitle{Experiences}{\faHistory}
        \begin{experiences}
        
        \experience
        {\DTMlangsetup*{showdayofmonth=false}Present}{Research Associate | Scene Analysis \& Activity Recognition}{Stuttgart}{Fraunhofer IPA}
        {\DTMlangsetup*{showdayofmonth=false}\DTMdisplaydate{2024}{02}{01}}
        {   \begin{itemize}
                % ORDER IS PER PROFILE (since 2026-09-22). Each bullet below is
                % defined once, in place, as \cvFhg<Name>, with its notes beside
                % it, and nothing prints until \cvFraunhoferOrder at the end of
                % this list. cv.tex sets the default order, which is the order
                % written here. A profile may \renewcommand it to lead with what
                % its posting asks for first, and terabase-cv does. Reordering
                % selects and orders; it does not fork the prose. Keep the
                % German branch's \def names identical.
                %
                % Ordered for this reader, not chronologically: the platform
                % cluster first -- CI/CD, builds, gateway, docs -- then the
                % research work. Re-ordered when \ifcvshowcudakernels went in:
                % hiding the kernels bullet on a platform profile silently
                % promoted the labeling result to the opening line, which is not
                % what a platform reader is looking for. The research profiles
                % still show the kernels bullet, so it still leads there.
                %
                % Was eight bullets; a reader gives an entry four or five before
                % skimming, and the last three were carrying real work into that
                % blind spot. Two merges did it, and both removed a repetition
                % rather than a fact:
                %   - the camera-pipeline bullet and the "computer vision
                %     workloads across GPU servers and edge devices" bullet were
                %     the same work seen twice, once from the client end and once
                %     from the hardware end. Joined, they show the whole span --
                %     which is the point, and which neither said alone.
                %   - CI/CD and DGX administration are one sentence: the platform
                %     the group runs on. The old CI/CD bullet also spent half a
                %     line explaining that CI/CD means build, test and deploy.
                %
                % "Replaced ... hand-labeling" and not "Automated data labeling":
                % same fact, but one names what it cost the group before. Verbs
                % that state an outcome are the only thing standing in for the
                % numbers this list still does not have.
                % The GLM-5.2 optimisations were a parent bullet plus two
                % sub-bullets: six lines for what fits in three as one sentence.
                % A nested itemize also indents twice, so each child line is
                % narrower and wraps sooner -- it costs lines twice over. Same
                % numbers, same order, no sub-list.
                \def\cvFhgKernels{\ifcvshowcudakernels
                \item Accelerated the team's efficient-learning and optimization methods with custom CUDA kernels and Python bindings to C++
                \fi}
                % Fraunhofer IPA's quantum-computing group. Added for the amd-hpc
                % profile (see the flag in cv.tex); ~700x is the owner's measured
                % figure. The link is the public project page, which is what makes
                % the claim checkable -- keep it. Replaces the RemaNet bullet on
                % that profile: edge-AI product work is the one item here an HPC
                % posting does not buy.
                \def\cvFhgQuantum{\ifcvshowquantum
                \item Ported a \textbf{quantum-physics simulation} from CPU to GPU for the \link{https://www.kqcbw.de}{KQCBW} with \textbf{CUDA} and \textbf{PyTorch}, profiling with \textbf{Nsight} and \textbf{VTune} to a \textbf{$\sim$700$\times$} speedup
                \fi}
                % Opens the platform cluster -- and on a platform profile, where
                % the kernels bullet above is hidden, opens the entry. It was the
                % only bullet here carrying no number, which for a platform
                % posting was the wrong one to leave bare. Ratio deliberately
                % unpinned: AGENTS.md forbids quoting a wall-clock digit that no
                % gate re-measures. Cache hit rate is measurable per run and is
                % the figure to swap in if one is ever wanted.
                % \ifcvshowmldetail clause (owner, 2026-09-22): the group
                % tracked experiments in MLflow and ClearML -- the posting's
                % "experiment tracking". Passive on purpose: these were the
                % team's tools, used, not owned like the CI/CD before them.
                \def\cvFhgCicd{%
                % "Docker" behind \ifcvshowstackdetail (owner, 2026-09-23: Docker
                % daily at Fraunhofer, advanced builds; BuildKit is its build
                % engine). Off, the sentence is exactly the line below. That
                % line, kept:
                % \item Owned the institute's GitLab CI/CD, DGX A100/H200 GPUs and \textbf{Grafana/Prometheus} monitoring; \textbf{BuildKit} and \textbf{sccache} caching cut builds from hours to minutes\ifcvshowmldetail; experiments tracked in \textbf{MLflow} and \textbf{ClearML}\fi
                \item Owned the institute's GitLab CI/CD, DGX A100/H200 GPUs and \textbf{Grafana/Prometheus} monitoring; \textbf{BuildKit} and \textbf{sccache} caching cut \ifcvshowstackdetail\textbf{Docker} \fi builds from hours to minutes\ifcvshowmldetail; experiments tracked in \textbf{MLflow} and \textbf{ClearML}\fi
                }
                % Naming the projects is the point: "fixes go upstream" is a
                % posture, named projects are checkable evidence of root-cause
                % work in somebody else's codebase -- the posting's "root cause
                % analysis" ask, answered nowhere else on the page.
                %
                % LLVM and sccache, and "upstreamed" rather than a PR count:
                % owner decision 2026-09-04. ANTfrastructure's
                % https://github.com/Kataglyphis/ANTfrastructure/blob/develop/docs/upstream-windows-patches.md "Already filed" table (checked
                % 2026-09-02) records LLVM #219275/#219276/#219200 and hcsshim
                % #2855 as OPEN, and lists sccache #2808 as an issue we opened
                % rather than a patch of ours -- but that register tracks only
                % what this repo filed, and it is the owner's own contribution
                % history, not a three-day-old table, that settles this.
                %
                % Two things follow. "Upstreamed" reads as MERGED, so it has to
                % be merged if a reviewer opens the PR list. And do not "restore"
                % hcsshim or a PR count here from the register: both were
                % deliberately removed.
                % Behind \ifcvshowbuilds (default on; see cv.tex).
                \def\cvFhgBuilds{\ifcvshowbuilds
                \item Advanced \textbf{CMake} build systems for C++ codebases, \textbf{GStreamer} (\textbf{Meson}) and \textbf{LiteRT-LM} (\textbf{Bazel}) built from source; fixes upstreamed to \textbf{LLVM} and \textbf{sccache}, not kept as local patches
                \fi}
                % Third in the platform cluster, and the only place the CV claims
                % Go -- see the flag in cv.tex. "45 researchers" is the scale
                % figure: the plugins are infrastructure a whole department runs
                % through, which a bare "token metering" does not convey.
                % "access-controlled" (owner, 2026-09-23): the gateway
                % authenticates its users, so it is access control, not only
                % metering -- the c12-automation posting's "private networking
                % and access control". It describes the gateway; the owner's
                % part stays the plugins. The line before, kept:
                % \item Extended the institute's \textbf{AI gateway} with custom \textbf{WebAssembly} plugins in \textbf{Go}: token metering and rate limiting for its \textbf{45 researchers}
                \def\cvFhgGateway{\ifcvshowgateway
                \item Extended the institute's access-controlled \textbf{AI gateway} with custom \textbf{WebAssembly} plugins in \textbf{Go}: token metering and rate limiting for its \textbf{45 researchers}
                \fi}
                % Closes the platform cluster, deliberately after the systems it
                % documents: the posting asks for documentation "that will help
                % them make the most out of the tools and systems you'll build",
                % which is this order exactly. It names those systems rather than
                % saying "docs" -- an abstract documentation claim is the easiest
                % thing on a CV for a reader to discount.
                % One line since 2026-09-23 (c12-automation, the first
                % profile to render it next to the others, set "from" alone on
                % a second line). Same facts. The first wording, kept:
                % \item Wrote the runbooks and onboarding docs for the CI/CD and GPU cluster the group works from
                % Owner-confirmed 2026-10-25; no repository count or review responsibilities inferred.
                \def\cvFhgMaintainer{\ifcvshowmaintainer
                \item Code maintainer for several institute repositories on \textbf{GitLab}, in \textbf{Python}, \textbf{Rust} and \textbf{C++}
                \fi}
                \def\cvFhgDocs{\ifcvshowdocs
                \item Wrote the runbooks and onboarding docs for the group's CI/CD and GPU cluster
                \fi}
                % "Replaced ... hand-labeling" and not "Automated data labeling":
                % same fact, but one names what it cost the group before. Verbs
                % that state an outcome are the only thing standing in for the
                % numbers this list still does not have.
                % \ifcvshowmldetail clause (owner, 2026-09-22): the pipelines
                % run in the owner's own research tool for self-labeling, and
                % its labels partly became the training data of the owner's
                % projects -- "data annotation tooling" and dataset curation.
                % Not claimed: bounding boxes or a human review step (not
                % asked, not answered).
                % Reworded 2026-09-23 by the owner: the point is that some of
                % the owner's projects used the generated labels, not that
                % "its labels partly trained my models". The first clause,
                % kept:
                % \ifcvshowmldetail{} in my own research tool for self-labeling; its labels partly trained my models\fi
                \def\cvFhgVlm{%
                \item Cut hand-labeling time by \textbf{75\%} with \textbf{vLLM-served vision-language model pipelines}\ifcvshowmldetail{} in my own research tool for self-labeling; some of my projects used its generated labels\fi
                }
                % Both GLM-5.2 bullets are the same box: 6 GPUs of one H200
                % system. It said "6x H100" here until now, which read as a
                % second, different deployment. The hardware is named once, in
                % this bullet; the one below says "there" rather than repeating
                % it. Keep it that way -- naming it twice invites the reader to
                % assume two clusters, which is what the mismatch already did.
                % One deployment, three optimisations -- not two bullets. It was
                % split as "Deployed vLLM serving for GLM-5.2 on 6 GPUs..." and
                % "Ran GLM-5.2 (744B MoE) there at NVFP4...", which named the
                % model twice, the machine twice, and used two verbs for one
                % deploy. Six lines for what fits in four, and the split invited
                % the reader to count two systems. Every number is still here.
                % One headline number per optimisation -- and since 2026-09-04
                % only that: "-46% latency" came out because the bullet had
                % drifted back to four figures, and 10x throughput already makes
                % the speed claim. Two numbers a reader carries away beat four
                % they skim. The supporting figures were dropped earlier for the
                % same reason, not because they are weak:
                % ~143 tok/s at 55k context, 65% draft acceptance, -60% worst-
                % case latency, "no loss on large prefills", 8/8 perf gates, and
                % "than BF16" on the memory figure. Eight numbers in one bullet
                % is a spec sheet; three is a claim a reader carries away, and
                % the rest is what an interview is for. All of it is one edit
                % away if the page ever grows.
                % The coding agent IS this deployment -- it runs on the GLM-5.2
                % instance described here, so listing it separately claimed two
                % pieces of work where there is one, and the reader had no way
                % to tell. Purpose first, engineering after: what the cluster is
                % for, then what was done to make it fast enough to be usable.
                % "Started ... on our own initiative ... and grew it into
                % production" carries three things a claim never could: the work
                % was not assigned, it was done with someone, and it was carried
                % all the way. The posting asks for "collaborative, low-ego" --
                % this is the evidence, and evidence is the only way to answer
                % that ask without contradicting it by asserting it.
                % The opening stays -- unassigned, with someone, carried to
                % production is what this bullet is really evidence for. The
                % technical tail was a spec sheet: three techniques, each with
                % its own parenthesised metrics, five numbers in all. Kept the
                % two custom pieces of vLLM work and the memory result; dropped
                % TP2xPP3, "median", and the throughput/latency split per
                % technique. All of it is in the git history if it should return.
                % "a production service the institute's researchers use" and not
                % "production": the platform posting asks for environments built
                % FOR researchers, and this bullet is the only place the CV can
                % name its users. Says nothing new -- an institute coding agent
                % in production has exactly those users. An adoption figure
                % would be stronger still; there is no measured one to cite.
                % Before "privacy-first" (2026-09-22), kept for reference:
                % \item Started the institute's coding agent with a colleague on our own initiative and grew it into a production service its researchers use: self-hosted \textbf{GLM-5.2 and GLM-5.3} on \textbf{6 H200 GPUs}, with custom \textbf{speculative decoding} and \textbf{scheduling} (\textbf{$\sim$10$\times$ throughput}) and \textbf{NVFP4 + FP8 KV cache} (\textbf{1M-token context})
                % "privacy-first": the owner's word (2026-09-22) for why it is
                % self-hosted -- AI-first engineering without sending code to a
                % third party. It fits in the last line's slack.
                % "GLM-5.3 with vision capabilities" instead of "GLM-5.2 and
                % GLM-5.3": the owner's wording, 2026-09-23. It names the model
                % the service runs now and says it sees as well as writes,
                % which ties the agent to the vision work above it. GLM-5.2
                % stays true history (see the notes above). The version with
                % both models, kept:
                % \item Started the institute's coding agent with a colleague on our own initiative and grew it into a privacy-first production service its researchers use: self-hosted \textbf{GLM-5.2 and GLM-5.3} on \textbf{6 H200 GPUs}, with custom \textbf{speculative decoding} and \textbf{scheduling} (\textbf{$\sim$10$\times$ throughput}) and \textbf{NVFP4 + FP8 KV cache} (\textbf{1M-token context})
                \def\cvFhgAgent{%
                % The stack-detail clause first read "(open weights from
                % Hugging Face)" and set "context)" alone on a fifth line.
                % Kept: \ifcvshowstackdetail{} (open weights from \textbf{Hugging Face})\fi{}
                % "served with vLLM" added 2026-09-23 (owner: vLLM serves GLM-5.3
                % with vision). The bullet named the model but never the
                % engine; vLLM sat only in the tags. The line before, kept:
                % \item Started the institute's coding agent with a colleague on our own initiative and grew it into a privacy-first production service its researchers use: self-hosted \textbf{GLM-5.3 with vision capabilities}\ifcvshowstackdetail{} (from \textbf{Hugging Face})\fi{} on \textbf{6 H200 GPUs}, with custom \textbf{speculative decoding} and \textbf{scheduling} (\textbf{$\sim$10$\times$ throughput}) and \textbf{NVFP4 + FP8 KV cache} (\textbf{1M-token context})
                \item Started the institute's coding agent with a colleague on our own initiative and grew it into a privacy-first production service its researchers use: self-hosted \textbf{GLM-5.3 with vision capabilities}\ifcvshowstackdetail{} (from \textbf{Hugging Face})\fi{} served with \textbf{vLLM} on \textbf{6 H200 GPUs}, with custom \textbf{speculative decoding} and \textbf{scheduling} (\textbf{$\sim$10$\times$ throughput}) and \textbf{NVFP4 + FP8 KV cache} (\textbf{1M-token context})
                }
                % RemaNet = Remanufacturing Network, Horizon Europe grant
                % 101138627, 25 partners, Jan 2024 - Dec 2026, Fraunhofer IPA a
                % partner -- which is why it lines up with the February 2024
                % start date. Named because it is checkable: a reader can look
                % the project up, which "AI camera pipelines" alone does not
                % allow. The grant number stays out; the programme is the signal.
                % "Raspberry Pi/Hailo edge devices" undersold this: the Hailo
                % went into a Berghof PLC, i.e. an industrial controller was
                % retrofitted with an AI accelerator. That is a fact a reader can
                % picture and few CVs can claim; "edge devices" is a category.
                % Paid for by dropping "low-latency, high-throughput" and "end to
                % end" -- self-assessed adjectives, and the list from GPU server
                % down to PLC and up to the clients already says "end to end".
                % The mistral-rse builds have ~0pt spare, so this had to come out
                % no longer than it went in; check both before growing it.
                % The frontend is RemaNet's, not the PLC's. It used to trail the
                % sentence -- "...a Berghof PLC I fitted with a Hailo
                % accelerator, with a React web frontend" -- where it attaches
                % to the nearest noun, which is the PLC, and the second "with"
                % made that reading worse. Naming it first and hanging the
                % inference under it reads as one stack, top down, and everything
                % in the list belongs to RemaNet because RemaNet precedes the colon.
                % Linked like KQCBW in the quantum bullet: both are research
                % consortia a reader can look up, so both get the link colour
                % rather than bold. remanet-project.eu is the official site; it
                % names grant 101138627 and lists Fraunhofer IPA as a partner
                % (checked 2026-09-22). Not bold: the colour already marks the
                % name, and bold plus colour would outweigh KQCBW.
                % The bold, unlinked version, kept for reference:
                % \item Shipped \textbf{AI camera pipelines} for \textbf{RemaNet} (Horizon Europe): a \textbf{React} web frontend over C++ inference on NVIDIA GPU servers and a \textbf{Berghof PLC I fitted with a Hailo accelerator}
                %
                % CORRECTED 2026-09-22, by the owner. The Berghof PLC was NOT part
                % of RemaNet; it was a separate project. So was a third one, an
                % industry contract. The owner ran the whole ML lifecycle (design,
                % training, evaluation, deployment) in all three:
                %   RemaNet  -- activity detection; React frontend over C++
                %               inference on NVIDIA GPU servers (owner, later
                %               the same day: ONNX inference, PyTorch training,
                %               NVIDIA GPUs)
                %   Berghof  -- object detection on an embedded PLC fitted with a
                %               Hailo AI accelerator; the embedded target with a
                %               special AI chip was the challenge
                %   industry -- person detection for safety: 4 cameras at once,
                %               their 4 inferences batched
                % The notes above ("everything in the list belongs to RemaNet
                % because RemaNet precedes the colon") describe the superseded
                % bullet below. Its wording was right; the attribution was wrong.
                % Each project now has its own bullet and its own flag:
                % \ifcvshowremanet, \ifcvshowberghof and \ifcvshowsafety (see
                % cv.tex).
                % Superseded 2026-09-22 (it attributes the Berghof PLC to RemaNet):
                % \item Shipped \textbf{AI camera pipelines} for \link{https://remanet-project.eu}{RemaNet} (Horizon Europe): a \textbf{React} web frontend over C++ inference on NVIDIA GPU servers and a \textbf{Berghof PLC I fitted with a Hailo accelerator}
                % Before "25 partners" (2026-09-22), kept for reference:
                % \item Shipped \textbf{activity detection} for \link{https://remanet-project.eu}{RemaNet} (Horizon Europe): camera pipelines with a \textbf{React} web frontend over C++ inference on NVIDIA GPU servers
                % "25 partners": CORDIS and remanet-project.eu both list 25
                % organisations from 8 countries (checked 2026-09-22). It is the
                % cheapest evidence on the page of delivering inside a
                % multi-organisation consortium, i.e. working with external
                % partners. It stays inside the second line's slack, so it adds
                % no line to any profile.
                % With "25 partners" but before the owner's reformulation
                % (2026-09-22), kept for reference:
                % \item Shipped \textbf{activity detection} for \link{https://remanet-project.eu}{RemaNet} (Horizon Europe, 25 partners): camera pipelines with a \textbf{React} web frontend over C++ inference on NVIDIA GPU servers
                % Reformulated by the owner, 2026-09-22: inference runs on ONNX
                % and training in PyTorch, both on NVIDIA GPUs. "C++ inference"
                % is not wrong, but it named the language where the owner names
                % the runtime and the training framework. C++ stays in the tag
                % row. \slash, not "/": a bare slash glues "inference/PyTorch"
                % into one unbreakable word, which TeX can only fix with a bad
                % line. The owner's text ended in ";", which is dropped: no
                % other bullet ends in punctuation.
                \def\cvFhgRemanet{\ifcvshowremanet
                % \ifcvshowstackdetail clause (owner, 2026-09-23): TypeScript and
                % Node.js were and are used in RemaNet's web part.
                \item Shipped \textbf{activity detection} for \link{https://remanet-project.eu}{RemaNet} (Horizon Europe, 25 partners): camera pipelines with a \textbf{React} web frontend\ifcvshowstackdetail{} (\textbf{TypeScript}, \textbf{Node.js})\fi{} and \textbf{ONNX} inference\slash\textbf{PyTorch} training on NVIDIA GPUs
                \fi}
                % Before the frame rate and the framework (2026-09-22), kept:
                % \item Ran \textbf{object detection} on an embedded \textbf{Berghof PLC I fitted with a Hailo AI accelerator}
                % "~14 fps" and "PyTorch-trained": both from the owner,
                % 2026-09-22. The frame rate is measured on the Hailo chip. The
                % models were trained in PyTorch and deployed to the embedded
                % target -- the embedded chip was the challenge of this
                % project, and the pair says so in two words. Two lines now,
                % not one; the profiles that show this bullet have the room
                % (mistral-rse and amd-hpc hide it).
                % \ifcvshowmldetail clause (owner, 2026-09-22): the models ran
                % INT8-quantized on the Hailo chip, and the owner iterated with
                % the industry customers on field feedback -- the posting's
                % "gathering field feedback ... closing the loop between model
                % performance and real-world outcomes", answered with the
                % project where it happened most.
                % The clause's first wording, kept: it set "feedback" alone on a
                % third line, one line for one word.
                % \ifcvshowmldetail; \textbf{INT8}-quantized and iterated with industry customers on field feedback\fi
                \def\cvFhgBerghof{\ifcvshowberghof
                \item Ran PyTorch-trained \textbf{object detection} at \textbf{$\sim$14 fps} on an embedded \textbf{Berghof PLC I fitted with a Hailo AI accelerator}\ifcvshowmldetail; \textbf{INT8}-quantized, iterated on customer field feedback\fi
                \fi}
                % Earlier wordings, 2026-09-22, kept for reference:
                % \item Built \textbf{person detection} for industrial safety: \textbf{4 cameras} in parallel, their inferences \textbf{batched}
                % \item Built \textbf{person detection} for industrial safety: \textbf{4 cameras} in parallel, inferences \textbf{batched}
                % "Delivered ... for an industrial customer": it was a custom
                % industry project, so the model went to a paying customer and
                % had to hold up there. That is the "building, shipping and
                % supporting production B2B applications" a posting may ask for,
                % and "Built" did not say it.
                \def\cvFhgSafety{\ifcvshowsafety
                \item Delivered \textbf{person detection} for an industrial customer's safety application: \textbf{4 cameras} in parallel, inferences \textbf{batched}
                \fi}
                % KI-PrOT (owner, 2026-09-22; ki-prot.de checked the same day):
                % "Optimierte Produktionsplanung für die Oberflächentechnik",
                % funded by the BMFTR, coordinated by Softec, with Fraunhofer
                % IPA developing the AI methods. The owner built the visual
                % classification of parts into drag-out classes
                % (Verschleppungsklassen), created and annotated its dataset,
                % and versioned it with DVC. The site's stated aim is to reduce
                % chemical consumption; that is a goal, not a measured result,
                % so the bullet says "to reduce". It is the CV's only
                % classification work and its only DVC. Behind \ifcvshowkiprot,
                % default off.
                \def\cvFhgKiprot{\ifcvshowkiprot
                \item Built \textbf{visual classification} of parts into drag-out classes for \link{https://ki-prot.de}{KI-PrOT} to reduce chemical use in surface treatment: dataset created, annotated and versioned with \textbf{DVC}
                \fi}
                % Everything above is only defined. This line prints it, in the
                % order cv.tex or the profile sets.
                \cvFraunhoferOrder
            \end{itemize}
        }
        % Tags follow the bullets above, strongest first for this reader.
        % Dropped: "Edge AI", a category where the bullet now names the actual
        % accelerator, and "DGX A100/H200", which the CI/CD bullet and the GLM
        % one already spell out -- a tag that restates the prose above it only
        % costs width. Added NVFP4/FP8 (the quantisation work had no tag at all),
        % React and Hailo. LoRA stays although the prose no longer says it: the
        % fine-tuning is real, and a tag is the cheapest place to keep a keyword.
        % One line is the budget -- this row wrapping costs the mistral-rse
        % builds their last few points of slack.
        % PyTorch added 2026-09-22: it is on this entry's page (the quantum
        % port) and several postings list it as a must-have, and the row
        % still fits on one line. The list before, kept for reference:
        % {Python, Go, C++, CUDA, CMake, Bazel, Meson, vLLM, LoRA, GitLab CI/CD}
        % uv added the same day (owner): Python package and environment
        % management with Astral's uv, used at Fraunhofer. It sits next to
        % Python on purpose -- it is Python tooling. The tag is the tool's
        % own name, "uv". "Astral uv" was tried first and wrapped
        % "GitLab CI/CD" onto a second tag line, which would have cost every
        % profile a line, mistral-rse its page. The lists before, kept:
        % {Python, Go, C++, CUDA, PyTorch, CMake, Bazel, Meson, vLLM, LoRA, GitLab CI/CD}
        % {Python, Astral uv, Go, C++, CUDA, PyTorch, CMake, Bazel, Meson, vLLM, LoRA, GitLab CI/CD}
        % Per profile since 2026-09-23: \cvFhgTags in cv.tex holds the list
        % below as its default, and a profile may \renewcommand it. The
        % literal list it replaced here, kept:
        % {Python, uv, Go, C++, CUDA, PyTorch, CMake, Bazel, Meson, vLLM, LoRA, GitLab CI/CD}
        {\cvFhgTags}
        
        \emptySeparator
        
        \experience
        {\DTMlangsetup*{showdayofmonth=false}\DTMdisplaydate{2023}{05}{01}}{Research assistant}{Karlsruhe}{CAS Software AG}
        {\DTMlangsetup*{showdayofmonth=false}\DTMdisplaydate{2022}{10}{15}}  
        {    \ifcvshowcasdetail
            \begin{itemize}
                % Was three lowercase sentence fragments, one of which
                % ("providing user guidance in mixed reality applications")
                % only restated the entry it sits in.
                % Was two bullets that both ended on "guidance": the ONNX one
                % and "Applied eye tracking and then-new ML approaches to the
                % guidance problem". Eye tracking was the only fact the second
                % one added, so it moved into the first. One line back in German,
                % where the pair wrapped; in English it reads tighter either way.
                % One line since 2026-09-22 (the terabase-cv owner wanted this
                % entry described again, and the page had room for one line,
                % not two). Same facts; "with eye tracking" also stops reading
                % as if the models had it. The two-line version, kept:
                % \item Trained and deployed ONNX models with eye tracking for real-time user guidance in mixed reality
                \item Trained and deployed ONNX models for real-time eye-tracking guidance in mixed reality
            \end{itemize}
            \fi
        }
        {Machine learning, ONNX, Mixed reality, HoloLens2, C\#, Python}
        
        \end{experiences}
    }{
    
    \sectionTitle{Berufserfahrung}{\faHistory}
    \begin{experiences}
    % german
    \experience
    {\DTMlangsetup*{showdayofmonth=false}heute}{Wissenschaftlicher Mitarbeiter | Szenenanalyse \& Aktivitätserkennung}{Stuttgart}{Fraunhofer IPA}
    {\DTMlangsetup*{showdayofmonth=false}\DTMdisplaydate{2024}{02}{01}}
    {   \begin{itemize}
            % Wording is tighter than the English throughout this list, not for
            % taste: German sets ~10% longer and this entry is what decides
            % whether the Sprachen/Hobbies/Publikationen row fits on page one.
            % Reihenfolge und Zusammenlegungen wie im englischen Zweig -- siehe
            % die Begruendung dort. Der Nominalstil bleibt: er ist im deutschen
            % Lebenslauf ueblich und hier ausserdem das, was in die Zeile passt.
            % Siehe die Anmerkung im englischen Zweig: keine Unterliste.
            % REIHENFOLGE JE PROFIL (seit 2026-09-22), siehe die Anmerkung im
            % englischen Zweig: jeder Punkt wird hier nur als \cvFhg<Name>
            % definiert, ausgegeben wird er erst von \cvFraunhoferOrder am Ende
            % der Liste. Die Namen müssen denen im englischen Zweig gleichen.
            \def\cvFhgKernels{\ifcvshowcudakernels
            \item \textbf{CUDA-Kernels} und \textbf{Python-Bindings} an C++, die die Lern- und Optimierungsverfahren des Teams beschleunigen
            \fi}
            % Siehe die Anmerkung im englischen Zweig: Quantencomputing-Gruppe
            % des Fraunhofer IPA, gemessener ~700x-Speedup, der Link ist die
            % öffentliche Projektseite. Ersetzt dort den RemaNet-Punkt.
            \def\cvFhgQuantum{\ifcvshowquantum
            \item \textbf{Quantenphysikalische Simulation} des \link{https://www.kqcbw.de}{KQCBW} von CPU auf GPU portiert (\textbf{CUDA}, \textbf{PyTorch}) und mit \textbf{Nsight} und \textbf{VTune} auf \textbf{$\sim$700$\times$} Speedup profiliert
            \fi}
            % OFFEN: eine Zeile bleibt gedehnt (badness 5667), STRICT_WARNINGS
            % schlägt fehl. \allowbreak am Schrägstrich ändert den Umbruch nicht;
            % "Monitoring mit Grafana/Prometheus" löst es, kostet aber die Seite.
            \def\cvFhgCicd{%
            \item Verantwortlich für GitLab-CI/CD, DGX-A100/H200-GPUs und \textbf{Grafana/Prometheus}-Monitoring; \textbf{BuildKit}- und \textbf{sccache}-Caching verkürzten \ifcvshowstackdetail\textbf{Docker}-\fi Builds von Stunden auf Minuten\ifcvshowmldetail; Experiment-Tracking mit \textbf{MLflow} und \textbf{ClearML}\fi
            }
            % Siehe die Anmerkung im englischen Zweig.
            % Hinter \ifcvshowbuilds (standardmäßig an; siehe cv.tex).
            \def\cvFhgBuilds{\ifcvshowbuilds
            \item Fortgeschrittene \textbf{CMake}-Builds für C++-Codebasen, \textbf{GStreamer} (\textbf{Meson}) und \textbf{LiteRT-LM} (\textbf{Bazel}) aus dem Quellcode gebaut; Fixes an \textbf{LLVM} und \textbf{sccache} upstream eingebracht, statt sie lokal zu patchen
            \fi}
            % Siehe die Anmerkung im englischen Zweig.
            % "zugriffsgeschützt" seit 2026-09-23, siehe die Anmerkung im
            % englischen Zweig. Die Zeile davor, zum Nachschlagen erhalten:
            % \item \textbf{KI-Gateway} des Instituts um eigene \textbf{WebAssembly}-Plugins in \textbf{Go} erweitert: Token-Metering und Rate-Limiting für \textbf{45 Forschende}
            \def\cvFhgGateway{\ifcvshowgateway
            \item Zugriffsgeschütztes \textbf{KI-Gateway} des Instituts um eigene \textbf{WebAssembly}-Plugins in \textbf{Go} erweitert: Token-Metering und Rate-Limiting für \textbf{45 Forschende}
            \fi}
            % Siehe die Anmerkung im englischen Zweig: nennt die Systeme.
            \def\cvFhgMaintainer{\ifcvshowmaintainer
            \item Code-Maintainer mehrerer Instituts-Repositories auf \textbf{GitLab} in \textbf{Python}, \textbf{Rust} und \textbf{C++}
            \fi}
            \def\cvFhgDocs{\ifcvshowdocs
            \item Runbooks und Onboarding-Dokumentation für die CI/CD und den GPU-Cluster geschrieben, mit denen die Gruppe arbeitet
            \fi}
            % Siehe die Anmerkung im englischen Zweig. Umformuliert 2026-09-23;
            % die erste Fassung des Zusatzes, zum Nachschlagen erhalten:
            % \ifcvshowmldetail{} in meinem eigenen Forschungstool für Self-Labeling; seine Labels flossen teils ins Training meiner Modelle\fi
            \def\cvFhgVlm{%
            \item Manuelles Labeling um \textbf{75\%} reduziert: \textbf{vLLM-gehostete Vision-Language-Pipelines}\ifcvshowmldetail{} in meinem eigenen Forschungstool für Self-Labeling; einige meiner Projekte nutzten die generierten Labels\fi
            }
            % Siehe die Anmerkung im englischen Zweig: eine Maschine, einmal benannt.
            % Siehe die Anmerkung im englischen Zweig: ein Deployment, drei Optimierungen.
            % Siehe die Anmerkung im englischen Zweig: eine Kennzahl je Optimierung.
            % Siehe die Anmerkung im englischen Zweig: ein Deployment, nicht zwei.
            % Siehe die Anmerkung im englischen Zweig.
            % Siehe die Anmerkung im englischen Zweig.
            % Siehe die Anmerkung im englischen Zweig.
            % Vor "datenschutzkonform" (2026-09-22), zum Nachschlagen erhalten:
            % \item Coding-Agent des Instituts auf Eigeninitiative mit einem Kollegen als Dienst für die Forschenden in Produktion gebracht: selbst gehostetes \textbf{GLM-5.2 und GLM-5.3} auf \textbf{6 H200-GPUs}, mit \textbf{Speculative Decoding} und \textbf{Scheduling} (\textbf{$\sim$10$\times$ Durchsatz}) und \textbf{NVFP4 + FP8-KV-Cache} (\textbf{1M-Token-Kontext})
            % "GLM-5.3 mit Vision-Fähigkeiten" statt "GLM-5.2 und GLM-5.3"
            % (2026-09-23), siehe die Anmerkung im englischen Zweig. Die
            % Fassung mit beiden Modellen, zum Nachschlagen erhalten:
            % \item Coding-Agent des Instituts auf Eigeninitiative mit einem Kollegen als datenschutzkonformen Dienst für die Forschenden in Produktion gebracht: selbst gehostetes \textbf{GLM-5.2 und GLM-5.3} auf \textbf{6 H200-GPUs}, mit \textbf{Speculative Decoding} und \textbf{Scheduling} (\textbf{$\sim$10$\times$ Durchsatz}) und \textbf{NVFP4 + FP8-KV-Cache} (\textbf{1M-Token-Kontext})
            \def\cvFhgAgent{%
            \item Coding-Agent des Instituts auf Eigeninitiative mit einem Kollegen als datenschutzkonformen Dienst für die Forschenden in Produktion gebracht: selbst gehostetes \textbf{GLM-5.3 mit Vision-Fähigkeiten}\ifcvshowstackdetail{} (offene Gewichte von \textbf{Hugging Face})\fi{} mit \textbf{vLLM} auf \textbf{6 H200-GPUs}, mit \textbf{Speculative Decoding} und \textbf{Scheduling} (\textbf{$\sim$10$\times$ Durchsatz}) und \textbf{NVFP4 + FP8-KV-Cache} (\textbf{1M-Token-Kontext})
            }
            % Kein "end to end" hier: im Deutschen ein Anglizismus, den die
            % Doppelpunkt-Liste ohnehin schon sagt. Die Geraete stehen in einer
            % Klammer statt als Bindestrich-Kompositum: "Raspberry-Pi-/Hailo-
            % Edge-Geraeten" ist eine einzige unumbrechbare Einheit, und mit ihr
            % in der Zeile fand TeX keinen Umbruch mehr -- die erste Zeile wurde
            % auf Wortabstaende von mehreren Millimetern gezogen (Underfull
            % \hbox, badness 10000, was STRICT_WARNINGS zu Recht ablehnt).
            % NVIDIA faellt weg; der DGX-Punkt darunter nennt den Hersteller.
            % Siehe die Anmerkung im englischen Zweig.
            % "darauf" gestrichen: der Punkt lief um sechs Zeichen über zwei
            % Zeilen und kostete eine dritte fuer "native Apps" allein.
            % Siehe die Anmerkung im englischen Zweig.
            % Siehe die Anmerkung im englischen Zweig: das Frontend gehört zu RemaNet, nicht zur SPS.
            % Siehe die Anmerkung im englischen Zweig: verlinkt wie KQCBW.
            % Die fette, unverlinkte Fassung, zum Nachschlagen erhalten:
            % \item \textbf{KI-Kamera-Pipelines} für \textbf{RemaNet} (Horizon Europe): \textbf{React}-Web-Frontend über C++-Inferenz auf GPU-Servern und einer \textbf{Berghof-SPS mit Hailo}
            % KORRIGIERT 2026-09-22 (siehe die Anmerkung im englischen Zweig): die
            % Berghof-SPS war ein eigenes Projekt, nicht Teil von RemaNet.
            % Ersetzt 2026-09-22 (ordnet die Berghof-SPS RemaNet zu):
            % \item \textbf{KI-Kamera-Pipelines} für \link{https://remanet-project.eu}{RemaNet} (Horizon Europe): \textbf{React}-Web-Frontend über C++-Inferenz auf GPU-Servern und einer \textbf{Berghof-SPS mit Hailo}
            % Vor "25 Partner" (2026-09-22), zum Nachschlagen erhalten:
            % \item \textbf{Aktivitätserkennung} für \link{https://remanet-project.eu}{RemaNet} (Horizon Europe): Kamera-Pipelines mit \textbf{React}-Web-Frontend über C++-Inferenz auf GPU-Servern
            % Mit "25 Partner", vor der Umformulierung des Inhabers (2026-09-22),
            % zum Nachschlagen erhalten:
            % \item \textbf{Aktivitätserkennung} für \link{https://remanet-project.eu}{RemaNet} (Horizon Europe, 25 Partner): Kamera-Pipelines mit \textbf{React}-Web-Frontend über C++-Inferenz auf GPU-Servern
            % Siehe die Anmerkung im englischen Zweig: ONNX-Inferenz und
            % PyTorch-Training auf NVIDIA-GPUs.
            \def\cvFhgRemanet{\ifcvshowremanet
            \item \textbf{Aktivitätserkennung} für \link{https://remanet-project.eu}{RemaNet} (Horizon Europe, 25 Partner): Kamera-Pipelines mit \textbf{React}-Web-Frontend\ifcvshowstackdetail{} (\textbf{TypeScript}, \textbf{Node.js})\fi, \textbf{ONNX}-Inferenz\slash\textbf{PyTorch}-Training auf NVIDIA-GPUs
            \fi}
            % Vor Bildrate und Framework (2026-09-22), zum Nachschlagen erhalten:
            % \item \textbf{Objekterkennung} auf einer eingebetteten \textbf{Berghof-SPS}, die ich mit einem \textbf{Hailo}-KI-Beschleuniger nachgerüstet habe
            \def\cvFhgBerghof{\ifcvshowberghof
            \item In PyTorch trainierte \textbf{Objekterkennung} mit \textbf{$\sim$14 fps} auf einer eingebetteten \textbf{Berghof-SPS}, die ich mit einem \textbf{Hailo}-KI-Beschleuniger nachgerüstet habe\ifcvshowmldetail; \textbf{INT8}-quantisiert und mit Industriekunden anhand von Feldfeedback iteriert\fi
            \fi}
            % Frühere Fassung, 2026-09-22, zum Nachschlagen erhalten:
            % \item \textbf{Personenerkennung} für industrielle Sicherheit: \textbf{4 Kameras} parallel, Inferenzen \textbf{gebatcht}
            \def\cvFhgSafety{\ifcvshowsafety
            \item \textbf{Personenerkennung} für die Sicherheitsanwendung eines Industriekunden geliefert: \textbf{4 Kameras} parallel, Inferenzen \textbf{gebatcht}
            \fi}
            % Siehe die Anmerkung im englischen Zweig: KI-PrOT, Ziel und nicht
            % Messergebnis, daher "um ... zu sparen".
            \def\cvFhgKiprot{\ifcvshowkiprot
            \item \textbf{Visuelle Klassifikation} von Bauteilen in Verschleppungsklassen für \link{https://ki-prot.de}{KI-PrOT}, um Chemikalien in der Oberflächentechnik zu sparen: Datensatz erstellt, annotiert und mit \textbf{DVC} versioniert
            \fi}
            % Bis hier ist alles nur definiert; diese Zeile gibt es aus, in der
            % Reihenfolge aus cv.tex oder dem Profil.
            \cvFraunhoferOrder
        \end{itemize}
    }
    % Siehe die Anmerkung im englischen Zweig (PyTorch und uv seit
    % 2026-09-22). Frühere Listen, zum Nachschlagen erhalten:
    % {Python, Go, C++, CUDA, CMake, Bazel, Meson, vLLM, LoRA, GitLab CI/CD}
    % {Python, Go, C++, CUDA, PyTorch, CMake, Bazel, Meson, vLLM, LoRA, GitLab CI/CD}
    % {Python, Astral uv, Go, C++, CUDA, PyTorch, CMake, Bazel, Meson, vLLM, LoRA, GitLab CI/CD}
    % Seit 2026-09-23 je Profil (\cvFhgTags in cv.tex), siehe die Anmerkung im
    % englischen Zweig. Die ersetzte Liste, zum Nachschlagen erhalten:
    % {Python, uv, Go, C++, CUDA, PyTorch, CMake, Bazel, Meson, vLLM, LoRA, GitLab CI/CD}
    {\cvFhgTags}
        
    \emptySeparator
    
    \experience
    {\DTMlangsetup*{showdayofmonth=false}\DTMdisplaydate{2023}{05}{01}}   {Research assistant}{Karlsruhe}{CAS Software AG}
    {\DTMlangsetup*{showdayofmonth=false}
    \DTMdisplaydate{2022}{10}{15}} 
    {
        \ifcvshowcasdetail
        \begin{itemize}
            % Siehe die Anmerkung im englischen Zweig.
            % Nominal statt verbal, weil kürzer: die verbale Fassung ("Training
            % und Deployment von ONNX-Modellen für ...") lief um vier Zeichen
            % über die Zeile und kostete eine zweite für das Wortende.
            % Siehe die Anmerkung im englischen Zweig.
            \item ONNX-Modelle und Eye Tracking für Echtzeit-Benutzerführung in Mixed Reality
        \end{itemize}
        \fi
    }
    {Machine learning, ONNX, Mixed reality, HoloLens2, C\#, Python}
    %\emptySeparator
    \end{experiences}
}
%
    % The renderers: the creative-technology evidence on record, so they take
    % the room the hidden bullets left.
    % Awesome Source CV LaTeX Template
%
% This template has been downloaded from:
% https://github.com/darwiin/awesome-neue-latex-cv
%
% Author:
% Christophe Roger
%
% Template license:
% CC BY-SA 4.0 (https://creativecommons.org/licenses/by-sa/4.0/)
%Section: Project
\IfLanguageName{english}{
    \sectionTitle{Projects}{\faLaptop}

    \begin{projects}

        \project
    	{Renderer}{2020-now}
    	{\github{GraphicsEngineVulkan} \github{GraphicEngine}}
        {This projects provide me with a solid Vulkan/OpenGL starting point for implementing modern established rendering techniques and getting quickly started in own research topics.}
    	{CMake, C, C++, Python, Vulkan, OpenGL}
    				
        \ifcvshowprojml
        \project
    	{Machine Learning Algorithms}{2020 - now}
    	{\github{MachineLearningAlgorithms}}
    	{This project provides me a solid Machine Learning(ML) starting point for implementing modern established ML algorithms and getting quickly started in own research topics.}
    	{\LaTeX, Python, R}
        \fi

        \ifcvshowprojmr
        \project
    	{User-Adaptive MR Guidance}{2022-2023}
    	{\github{Designing-User-adaptive-Content-for-Mixed-Reality-Using-Eye-and-Hand-Tracking}}
    	{M.Sc.\ thesis: Investigates interrelationships between users' fixated objects, eye/head/hand movement, and intentions. Provides user-adaptive guidance for improved interaction quality in Mixed Reality.}
    	{\LaTeX, C\sharp, Mixed Reality, Python, Unity}
        \fi

        \ifcvshowprojbluenoise
        \project
    	{Blue Noise Error Distribution}{2019-2020}
    	{\github{BachelorArbeit}}
    	{B.Sc.\ thesis: Temporally stable blue noise error distribution in screen space for real-time rendering applications, building on prior work on blue noise error distributions for image quality.}
    	{\LaTeX, C, C++, Python}
        \fi

        \ifcvshowprojcollision
        \project
    	{GPU Collision Detection}{2021-2022}
    	{\github{GPUCollisionDetectionWithAS}}
    	{Novel collision detection using GPU acceleration structures. Implementation of a particle system accelerated on the GPU with BVH-based broad phase.}
    	{\LaTeX, C, C++, Vulkan API, GPGPU}
        \fi

    \end{projects}
}{
    \sectionTitle{Projekte}{\faLaptop}

    \begin{projects}
        \project
    	{Renderer}{2020-now}
    	{\github{GraphicsEngineVulkan}\github{GraphicEngine}}
    	{Diese Projekte bieten mir einen soliden Startpunkt um diverse etablierte Computergrafik Algorithmen nachzuimplementieren und darauf aufbauend auch eigene Forschung zu betreiben.}
    	{\LaTeX, CMake, C, C++, Python, Vulkan, OpenGL}
        
        \ifcvshowprojml
     	\project
     	{Machine Learning Algorithms}{2020 - now}
     	{\github{MachineLearningAlgorithms}}
     	{Dieses Projekt bietet mir einen soliden Startpunkt um diverse etablierte Machine Learning(ML) Algorithmen nachzuimplementieren und darauf aufbauend auch eigene Forschung zu betreiben.}
     	{\LaTeX, Python, R}
        \fi

        \ifcvshowprojmr
        \project
     	{Benutzeradaptive MR-Führung}{2022-2023}
     	{\github{Designing-User-adaptive-Content-for-Mixed-Reality-Using-Eye-and-Hand-Tracking}}
     	{Masterarbeit: Untersucht Wechselbeziehungen zwischen fixierten Objekten, Augen-/Kopf-/Handbewegungen und Absichten der Benutzer. Bietet benutzeradaptive Führung für verbesserte Interaktion in Mixed Reality.}
     	{\LaTeX, C\sharp, Mixed Reality, Python, Unity}
        \fi

        \ifcvshowprojbluenoise
        \project
     	{Blue Noise Fehlerverteilung}{2019-2020}
     	{\github{BachelorArbeit}}
     	{Bachelorarbeit: Zeitlich stabile Blue Noise Fehlerverteilung im Bildraum für Echtzeitanwendungen, aufbauend auf bisherigen Arbeiten zur Bildqualitätssteigerung.}
     	{\LaTeX, C, C++, Python}
        \fi

        \ifcvshowprojcollision
        \project
     	{GPU-Kollisionsdetektion}{2021-2022}
     	{\github{GPUCollisionDetectionWithAS}}
     	{Neuartige Kollisionsdetektion mit GPU-Beschleunigungsstrukturen. Implementierung eines GPU-beschleunigten Partikelsystems mit BVH-basierter Grobphase.}
     	{\LaTeX, C, C++, Vulkan API, GPGPU}
        \fi
        
    \end{projects}
}
%
    % YAAC Another Awesome CV LaTeX Template
%
% This template has been downloaded from:
% https://github.com/darwiin/yaac-another-awesome-cv
%
% Author:
% Christophe Roger
%
% Template license:
% CC BY-SA 4.0 (https://creativecommons.org/licenses/by-sa/4.0/)

%Section: Scholarships and additional info
\IfLanguageName{english}{
    
    \sectionTitle{Education}{\faMortarBoard}
    % Owner-confirmed 2026-10-25; coursework does not change the degree titles.
    \ifcvshowmedicalcoursework
    \noindent Medical imaging: electrical-engineering coursework during both bachelor's and master's studies.\par
    \fi
    \begin{experiences}
    \education
    {\DTMlangsetup*{showdayofmonth=false}
    \DTMdisplaydate{2023}{09}{30}}   {Master of Science Computer Science}{Karlsruhe}{Karlsruhe Institute of Technology}
    {\DTMlangsetup*{showdayofmonth=false}
    \DTMdisplaydate{2020}{04}{01}}
    % "1.9/ US-GPA: 3.1" had the slash glued to the number, which reads as an
    % unfinished fraction. More importantly the scale was missing: German grades
    % run 1.0 (best) to 4.0 (pass), so a reader used to higher-is-better -- an
    % American, or a French one on the 0-20 scale -- reads 1.9 as a poor result.
    % The US-GPA beside it only half fixes that; naming the scale fixes it.
    % The German branch keeps the bare "Gesamtnote: 1,9" on purpose: for a German
    % reader the scale is not in question, and the US-GPA would be noise.
    {Overall grade: 1.9 (German scale, 1.0 = best), US-GPA: 3.1}
    {
        % Thesis lines are the first education detail a profile can afford to
        % lose: the degree, its dates, its grade and its subject tags all stay,
        % and a thesis title is rarely what an industry reader is deciding on.
        % Dropping whole entries instead would leave visible holes in the dates.
        \ifcvshowthesislines
        \begin{itemize}
            \item Thesis: Designing User-adaptive Guidance for Mixed Reality Using Eye Tracking
        \end{itemize}
        \fi
    }
    % Python added 2026-09-22 (owner). The list before, kept for reference:
    % {AI, Computer Graphics and Geometry Processing, Anthropomatics and Cognitive Systems}
    {AI, Computer Graphics and Geometry Processing, Anthropomatics and Cognitive Systems, Python}
    \emptySeparator
    \education{\DTMlangsetup*{showdayofmonth=false}
    \DTMdisplaydate{2020}{03}{01}}
    {Bachelor of Science Computer Science}
    {Karlsruhe}
    {Karlsruhe Institute of Technology}
    { \DTMlangsetup*{showdayofmonth=false}
    \DTMdisplaydate{2015}{10}{01}}
    {Overall grade: 2.5 (German scale, 1.0 = best), US-GPA: 2.5}
    {
    \ifcvshowthesislines
    \begin{itemize}
        \item Thesis: Temporally Stable Blue Noise Error Distribution in Screen Space for Real Time Applications
    \end{itemize}
    \fi
    }
    {Path Tracing, Blue Noise, C++, Real-Time Systems, Low Latency, High Performance}
    \ifcvshowschool
    \emptySeparator
    \education{2014}
    {Pupil | Higher education entrance qualification (A-levels)}
    {Schwäbisch Hall}
    {Gymnasium bei St. Michael}
    {2006}
    {Overall grade: 2.2 (German scale, 1.0 = best), US-GPA: 2.8}
    {
        % Commented out to fit the CV on one page -- restore by uncommenting.
        %\begin{itemize}
        %    \item First steps in Procedural programming: Sudoku project
        %    \item Started learning Java autodidactically
        %\end{itemize}
    }
    {Delphi, Pascal, Java}
    \fi
    \end{experiences}

}{
    
    \sectionTitle{Bildung}{\faMortarBoard}
    \ifcvshowmedicalcoursework
    \noindent Medizinische Bildgebung: Elektrotechnik-Lehrveranstaltungen im Bachelor- und Masterstudium.\par
    \fi
    \begin{experiences}
    % german
    \education
    {\DTMlangsetup*{showdayofmonth=false}
    \DTMdisplaydate{2023}{09}{30}}   {Master of Science Informatik}{Karlsruhe}{Karlsruher Institut für Technologie}
    {\DTMlangsetup*{showdayofmonth=false}
    \DTMdisplaydate{2020}{04}{01}}
    {Gesamtnote: 1,9}
    {
        % Siehe die Anmerkung im englischen Zweig.
        \ifcvshowthesislines
        \begin{itemize}
            % "Thesis:" ist im deutschen Lebenslauf ein Anglizismus -- die
            % Arbeiten heißen Master- und Bachelorarbeit. Das englische
            % "Thesis:" drüben bleibt, es ist dort das richtige Wort.
            \item Masterarbeit: Entwurf von benutzeranpassbarem Szeneninhalt für Mixed Reality mit Hilfe von Augen- und Handnachverfolgung
        \end{itemize}
        \fi
    }
    % KI, nicht AI: die deutschen Aufzählungspunkte schreiben durchgehend
    % KI-Kamera-Pipelines und Hailo-KI-Beschleuniger.
    % Python seit 2026-09-22. Die Liste davor, zum Nachschlagen erhalten:
    % {KI, Computergrafik und Geometrieverarbeitung, Anthropomatik und Kognitive Systeme}
    {KI, Computergrafik und Geometrieverarbeitung, Anthropomatik und Kognitive Systeme, Python}
    \emptySeparator
    \education{\DTMlangsetup*{showdayofmonth=false}
    \DTMdisplaydate{2020}{03}{01}}
    {Bachelor of Science Informatik}
    {Karlsruhe}
    {Karlsruher Institut für Technologie}
    { \DTMlangsetup*{showdayofmonth=false}
    \DTMdisplaydate{2015}{10}{01}}
    {Gesamtnote: 2,5}
    {
    % Siehe die Anmerkung im englischen Zweig.
    \ifcvshowthesislines
    \begin{itemize}
        % Siehe die Anmerkung beim Master.
        \item Bachelorarbeit: Zeitlich stabile Blue Noise Fehlerverteilung im Bildraum für Echtzeitanwendungen
    \end{itemize}
    \fi
    }
    {Path Tracing, Blue Noise, C++, Echtzeitsysteme, Low Latency, High Performance}
    \ifcvshowschool
    \emptySeparator
    \education{2014}
    {Schüler | Allgemeine Hochschulreife}
    {Schwäbisch Hall}
    {Gymnasium bei St. Michael}
    {2006}
    {Gesamtnote: 2,2}
    {
        % Commented out to fit the CV on one page -- restore by uncommenting.
        %\begin{itemize}
        %    \item erste Schritte in proceduraler Programmierung: Sudokuprojekt
        %    \item Java autodidaktisch gelernt
        %\end{itemize}
    }
    {Delphi, Pascal, Java}
    \fi
    \end{experiences}
}
  


%
    % YAAC Another Awesome CV LaTeX Template
%
% This template has been downloaded from:
% https://github.com/darwiin/yaac-another-awesome-cv
%
% Author:
% Christophe Roger
%
% Template license:
% CC BY-SA 4.0 (https://creativecommons.org/licenses/by-sa/4.0/)

% Languages / Hobbies / Publications share one \threecolumnsection row.
% It is the last thing on a strict one-pager, so the room it gets here (full
% size, 1em under each title, a gap above) is paid for elsewhere: the trimmed
% \LTpre/\LTpost around the experiences longtables in myCV_METADATA.cls, and
% the smaller post-header \vspace in cv.tex. Grow any of those and this row
% stops fitting -- rebuild and check the page count.
\IfLanguageName{english}{
    %Section: Languages
    \vspace{0.05cm}
    \par\noindent
    \threecolumnsection
    {\sectionTitle{Languages}{\faLanguage}
    \vspace{0.55em}
    \begin{itemize}
    	\item German (native)
    	\item English, level C1
    	\item French, level B1
    	\item Latin
    \end{itemize}}
    {\sectionTitle{Hobbies}{\faBicycle}
    \vspace{0.55em}
    \begin{itemize}
        % Open source leads this list on purpose: it is the one hobby the
        % rest of the page is evidence for. Costs nothing -- same item,
        % same column height, first line instead of second.
        \item Open source
        % Powerlifting out, Cycling in. The two branches disagreed here anyway
        % -- English said Powerlifting, German said Powerbuilding -- so that
        % mismatch goes with it. One in, one out: the column keeps its five
        % lines, and this row's height stays where it was.
        \item Cycling
        % Commented out per request -- restore by uncommenting.
        % \item Bouldering
        \item Guitar
        \item Philosophy
    \end{itemize}
    }
    {\sectionTitle{Publications}{\faFileTextO}
    % \raggedright as in \twocolumnsection: justifying one long title across
    % this wide column stretches the interword spaces visibly.
    \raggedright
    \vspace{0.55em}
    \begin{itemize}
        \item \textquote{Semantic Robustness Probing via Inpainting} -- ECML PKDD 2026 Demo Track (with N.\ Steckhan)
        % The SemProbe one-liner is commented out, not deleted: it cost two
        % \footnotesize lines, and with a second publication in the column
        % Publications is now the tallest of the three and sets the row height
        % on its own. An abstract is worth less on a CV than a second named,
        % lookup-able deliverable. Restore it by uncommenting -- and then check
        % the German mistral-rse build, which is the one it pushes over.
        % \item[] {\footnotesize \textsc{SemProbe}: diffusion-based inpainting probes object detectors against ODD-derived factors.}
        % The named study is the concrete deliverable a reader can look up, and
        % it is Fraunhofer work rather than a side project -- it benchmarks
        % accelerator hardware for computer vision, which is the evidence behind
        % the "GPU acceleration" and edge-device claims above. German title kept
        % as published; the gloss after the dash is what an English reader needs.
        % "co-authored" rather than a bare title: the ECML entry above names its
        % co-author, so a second publication with no such mark reads as solo work.
        % Swap in the name when it is to hand -- matching the entry above is better.
        \item \textquote{KI-Beschleunigung durch spezialisierte Hardware und Software} -- co-authored study on AI accelerators, Fraunhofer IPA/IAO (2026)
    \end{itemize}
    }
}{
    %Section: Languages
    \vspace{0.05cm}
    \par\noindent
    \threecolumnsection
    {\sectionTitle{Sprachen}{\faLanguage}
    \vspace{0.55em}
    \begin{itemize}
    	\item Deutsch (Muttersprache)
    	% CEFR levels in parentheses: ", Level B1" overflows this narrow column.
    	\item Englisch (C1)
    	\item Französisch (B1)
    	\item Latein
    \end{itemize}}
    {\sectionTitle{Hobbies}{\faBicycle}
    \vspace{0.55em}
    \begin{itemize}
        % Siehe die Anmerkung im englischen Zweig.
        \item Open Source
        \item Radfahren
        \item Gitarre
        \item Philosophie
    \end{itemize}
    }
    {\sectionTitle{Publikationen}{\faFileTextO}
    % \raggedright as in \twocolumnsection: justifying one long title across
    % this wide column stretches the interword spaces visibly.
    \raggedright
    \vspace{0.55em}
    \begin{itemize}
        \item \textquote{Semantic Robustness Probing via Inpainting} -- ECML PKDD 2026 Demo Track (mit N.\ Steckhan)
        % Siehe die Anmerkung im englischen Zweig.
        % \item[] {\footnotesize \textsc{SemProbe}: Diffusions-Inpainting prüft Objektdetektoren gegen ODD-Faktoren.}
        % Siehe die Anmerkung im englischen Zweig.
        % Siehe die Anmerkung im englischen Zweig.
        \item \textquote{KI-Beschleunigung durch spezialisierte Hardware und Software} -- Studie, Fraunhofer IPA/IAO (2026)
    \end{itemize}
    }
}
%
}
